% fig-architecture.tex — 圖 1：四層偵測架構（繁體中文版）
% 自包含的 TikZ 圖。以 xelatex 編譯。
% 由 main.tex 透過 % fig-architecture.tex — Figure 1: Four-layer detection architecture
% Self-contained TikZ figure. Compile with pdflatex/xelatex/lualatex.
% Used by main.tex via % fig-architecture.tex — Figure 1: Four-layer detection architecture
% Self-contained TikZ figure. Compile with pdflatex/xelatex/lualatex.
% Used by main.tex via % fig-architecture.tex — Figure 1: Four-layer detection architecture
% Self-contained TikZ figure. Compile with pdflatex/xelatex/lualatex.
% Used by main.tex via \input{fig-architecture}.

\begin{figure}[t]
\centering
\begin{tikzpicture}[
  node distance=4mm and 6mm,
  every node/.style={font=\small},
  layer/.style={
    draw, rounded corners=2pt,
    minimum width=42mm, minimum height=10mm,
    align=center, fill=gray!8,
  },
  surface/.style={layer, fill=blue!7},
  behavior/.style={layer, fill=blue!12},
  v3/.style={layer, fill=orange!12},
  v3e/.style={layer, fill=orange!18},
  fusion/.style={layer, fill=green!12, minimum width=88mm, minimum height=12mm},
  input/.style={
    draw, rounded corners=2pt,
    minimum width=88mm, minimum height=8mm,
    align=center, fill=gray!4,
  },
  output/.style={
    draw, rounded corners=2pt,
    minimum width=88mm, minimum height=8mm,
    align=center, fill=gray!4,
    font=\small\itshape,
  },
  arrow/.style={-{Latex[length=1.6mm]}, thick, gray!70!black},
  label/.style={font=\scriptsize\itshape, gray!50!black},
]

% Top: input
\node[input] (input) {Probe responses from \texttt{(baseUrl, model, key)}};

% Layer 1 + Layer 2 side by side
\node[surface, below left=10mm and 0mm of input] (l1) {%
  \textbf{Layer \textcircled{\scriptsize 1}}\\Surface fingerprint\\(family weights)};
\node[behavior, below right=10mm and 0mm of input] (l2) {%
  \textbf{Layer \textcircled{\scriptsize 2}}\\Behavior column\\(self-claim zeroed)};

% Layer 3 below
\node[v3, below=10mm of input, yshift=-22mm] (l3) {%
  \textbf{Layer \textcircled{\scriptsize 3}} \quad Sub-model deterministic match\\
  cutoff $\,\cdot\,$ capability $\,\cdot\,$ refusal-prefix};
\path (l3.west) -- node[layer, minimum width=88mm, draw=none, fill=none] {} (l3.east);

% Layer 4 below
\node[v3e, below=4mm of l3] (l4) {%
  \textbf{Layer \textcircled{\scriptsize 4}} \quad Behavioral-vector extension\\
  refusal ladder (8) $\,\cdot\,$ formatting (3) $\,\cdot\,$ uncertainty (1)};

% Verdict fusion
\node[fusion, below=4mm of l4] (fuse) {%
  \textbf{Verdict fusion}\\
  7 mutually exclusive states + confidence band};

% Output
\node[output, below=4mm of fuse] (out) {%
  IdentityAssessment \{ verdict, riskFlags, evidence, confidence \}};

% Arrows
\draw[arrow] (input.south) -- ++(0,-2mm) -| (l1.north);
\draw[arrow] (input.south) -- ++(0,-2mm) -| (l2.north);
\draw[arrow] (l1.south) -- ++(0,-2mm) -| ([xshift=-22mm]l3.north);
\draw[arrow] (l2.south) -- ++(0,-2mm) -| ([xshift=22mm]l3.north);
\draw[arrow] (l3.south) -- (l4.north);
\draw[arrow] (l4.south) -- (fuse.north);
\draw[arrow] (fuse.south) -- (out.north);

% Side annotations
\node[label, right=2mm of l1, anchor=west, text width=20mm] {polluted by self-claim};
\node[label, left=2mm of l2, anchor=east, text width=20mm] {robust to prompt-layer attack};
\node[label, right=2mm of l3, anchor=west, text width=20mm] {can be precisely targeted};
\node[label, right=2mm of l4, anchor=west, text width=20mm] {policy-baked, hard to forge};

\end{tikzpicture}
\caption{The four-layer independent-fingerprint detection architecture (\S\ref{sec:method}).
Each layer uses an independent signal channel; an adversary must simultaneously forge
consistent behavior across multiple channels to escape detection. The verdict fusion
emits one of seven mutually exclusive states with a confidence band.}
\label{fig:architecture}
\end{figure}
.

\begin{figure}[t]
\centering
\begin{tikzpicture}[
  node distance=4mm and 6mm,
  every node/.style={font=\small},
  layer/.style={
    draw, rounded corners=2pt,
    minimum width=42mm, minimum height=10mm,
    align=center, fill=gray!8,
  },
  surface/.style={layer, fill=blue!7},
  behavior/.style={layer, fill=blue!12},
  v3/.style={layer, fill=orange!12},
  v3e/.style={layer, fill=orange!18},
  fusion/.style={layer, fill=green!12, minimum width=88mm, minimum height=12mm},
  input/.style={
    draw, rounded corners=2pt,
    minimum width=88mm, minimum height=8mm,
    align=center, fill=gray!4,
  },
  output/.style={
    draw, rounded corners=2pt,
    minimum width=88mm, minimum height=8mm,
    align=center, fill=gray!4,
    font=\small\itshape,
  },
  arrow/.style={-{Latex[length=1.6mm]}, thick, gray!70!black},
  label/.style={font=\scriptsize\itshape, gray!50!black},
]

% Top: input
\node[input] (input) {Probe responses from \texttt{(baseUrl, model, key)}};

% Layer 1 + Layer 2 side by side
\node[surface, below left=10mm and 0mm of input] (l1) {%
  \textbf{Layer \textcircled{\scriptsize 1}}\\Surface fingerprint\\(family weights)};
\node[behavior, below right=10mm and 0mm of input] (l2) {%
  \textbf{Layer \textcircled{\scriptsize 2}}\\Behavior column\\(self-claim zeroed)};

% Layer 3 below
\node[v3, below=10mm of input, yshift=-22mm] (l3) {%
  \textbf{Layer \textcircled{\scriptsize 3}} \quad Sub-model deterministic match\\
  cutoff $\,\cdot\,$ capability $\,\cdot\,$ refusal-prefix};
\path (l3.west) -- node[layer, minimum width=88mm, draw=none, fill=none] {} (l3.east);

% Layer 4 below
\node[v3e, below=4mm of l3] (l4) {%
  \textbf{Layer \textcircled{\scriptsize 4}} \quad Behavioral-vector extension\\
  refusal ladder (8) $\,\cdot\,$ formatting (3) $\,\cdot\,$ uncertainty (1)};

% Verdict fusion
\node[fusion, below=4mm of l4] (fuse) {%
  \textbf{Verdict fusion}\\
  7 mutually exclusive states + confidence band};

% Output
\node[output, below=4mm of fuse] (out) {%
  IdentityAssessment \{ verdict, riskFlags, evidence, confidence \}};

% Arrows
\draw[arrow] (input.south) -- ++(0,-2mm) -| (l1.north);
\draw[arrow] (input.south) -- ++(0,-2mm) -| (l2.north);
\draw[arrow] (l1.south) -- ++(0,-2mm) -| ([xshift=-22mm]l3.north);
\draw[arrow] (l2.south) -- ++(0,-2mm) -| ([xshift=22mm]l3.north);
\draw[arrow] (l3.south) -- (l4.north);
\draw[arrow] (l4.south) -- (fuse.north);
\draw[arrow] (fuse.south) -- (out.north);

% Side annotations
\node[label, right=2mm of l1, anchor=west, text width=20mm] {polluted by self-claim};
\node[label, left=2mm of l2, anchor=east, text width=20mm] {robust to prompt-layer attack};
\node[label, right=2mm of l3, anchor=west, text width=20mm] {can be precisely targeted};
\node[label, right=2mm of l4, anchor=west, text width=20mm] {policy-baked, hard to forge};

\end{tikzpicture}
\caption{The four-layer independent-fingerprint detection architecture (\S\ref{sec:method}).
Each layer uses an independent signal channel; an adversary must simultaneously forge
consistent behavior across multiple channels to escape detection. The verdict fusion
emits one of seven mutually exclusive states with a confidence band.}
\label{fig:architecture}
\end{figure}
.

\begin{figure}[t]
\centering
\begin{tikzpicture}[
  node distance=4mm and 6mm,
  every node/.style={font=\small},
  layer/.style={
    draw, rounded corners=2pt,
    minimum width=42mm, minimum height=10mm,
    align=center, fill=gray!8,
  },
  surface/.style={layer, fill=blue!7},
  behavior/.style={layer, fill=blue!12},
  v3/.style={layer, fill=orange!12},
  v3e/.style={layer, fill=orange!18},
  fusion/.style={layer, fill=green!12, minimum width=88mm, minimum height=12mm},
  input/.style={
    draw, rounded corners=2pt,
    minimum width=88mm, minimum height=8mm,
    align=center, fill=gray!4,
  },
  output/.style={
    draw, rounded corners=2pt,
    minimum width=88mm, minimum height=8mm,
    align=center, fill=gray!4,
    font=\small\itshape,
  },
  arrow/.style={-{Latex[length=1.6mm]}, thick, gray!70!black},
  label/.style={font=\scriptsize\itshape, gray!50!black},
]

% Top: input
\node[input] (input) {Probe responses from \texttt{(baseUrl, model, key)}};

% Layer 1 + Layer 2 side by side
\node[surface, below left=10mm and 0mm of input] (l1) {%
  \textbf{Layer \textcircled{\scriptsize 1}}\\Surface fingerprint\\(family weights)};
\node[behavior, below right=10mm and 0mm of input] (l2) {%
  \textbf{Layer \textcircled{\scriptsize 2}}\\Behavior column\\(self-claim zeroed)};

% Layer 3 below
\node[v3, below=10mm of input, yshift=-22mm] (l3) {%
  \textbf{Layer \textcircled{\scriptsize 3}} \quad Sub-model deterministic match\\
  cutoff $\,\cdot\,$ capability $\,\cdot\,$ refusal-prefix};
\path (l3.west) -- node[layer, minimum width=88mm, draw=none, fill=none] {} (l3.east);

% Layer 4 below
\node[v3e, below=4mm of l3] (l4) {%
  \textbf{Layer \textcircled{\scriptsize 4}} \quad Behavioral-vector extension\\
  refusal ladder (8) $\,\cdot\,$ formatting (3) $\,\cdot\,$ uncertainty (1)};

% Verdict fusion
\node[fusion, below=4mm of l4] (fuse) {%
  \textbf{Verdict fusion}\\
  7 mutually exclusive states + confidence band};

% Output
\node[output, below=4mm of fuse] (out) {%
  IdentityAssessment \{ verdict, riskFlags, evidence, confidence \}};

% Arrows
\draw[arrow] (input.south) -- ++(0,-2mm) -| (l1.north);
\draw[arrow] (input.south) -- ++(0,-2mm) -| (l2.north);
\draw[arrow] (l1.south) -- ++(0,-2mm) -| ([xshift=-22mm]l3.north);
\draw[arrow] (l2.south) -- ++(0,-2mm) -| ([xshift=22mm]l3.north);
\draw[arrow] (l3.south) -- (l4.north);
\draw[arrow] (l4.south) -- (fuse.north);
\draw[arrow] (fuse.south) -- (out.north);

% Side annotations
\node[label, right=2mm of l1, anchor=west, text width=20mm] {polluted by self-claim};
\node[label, left=2mm of l2, anchor=east, text width=20mm] {robust to prompt-layer attack};
\node[label, right=2mm of l3, anchor=west, text width=20mm] {can be precisely targeted};
\node[label, right=2mm of l4, anchor=west, text width=20mm] {policy-baked, hard to forge};

\end{tikzpicture}
\caption{The four-layer independent-fingerprint detection architecture (\S\ref{sec:method}).
Each layer uses an independent signal channel; an adversary must simultaneously forge
consistent behavior across multiple channels to escape detection. The verdict fusion
emits one of seven mutually exclusive states with a confidence band.}
\label{fig:architecture}
\end{figure}
 引入。

\begin{figure*}[t]
\centering
\resizebox{\textwidth}{!}{%
\begin{tikzpicture}[
  node distance=4mm and 6mm,
  every node/.style={font=\small},
  layer/.style={
    draw, rounded corners=2pt,
    minimum width=42mm, minimum height=10mm,
    align=center, fill=gray!8,
  },
  surface/.style={layer, fill=blue!7},
  behavior/.style={layer, fill=blue!12},
  v3/.style={layer, fill=orange!12},
  v3e/.style={layer, fill=orange!18},
  fusion/.style={layer, fill=green!12, minimum width=88mm, minimum height=12mm},
  input/.style={
    draw, rounded corners=2pt,
    minimum width=88mm, minimum height=8mm,
    align=center, fill=gray!4,
  },
  output/.style={
    draw, rounded corners=2pt,
    minimum width=88mm, minimum height=8mm,
    align=center, fill=gray!4,
    font=\small\itshape,
  },
  arrow/.style={-{Latex[length=1.6mm]}, thick, gray!70!black},
  label/.style={font=\scriptsize\itshape, gray!50!black},
]

% 上方：輸入
\node[input] (input) {來自 \texttt{(baseUrl, model, key)} 的探測響應};

% 層 1 + 層 2 並排
\node[surface, below left=10mm and 0mm of input] (l1) {%
  \textbf{層 \textcircled{\scriptsize 1}}\\表面指紋\\（族別權重）};
\node[behavior, below right=10mm and 0mm of input] (l2) {%
  \textbf{層 \textcircled{\scriptsize 2}}\\行為對照\\（自我宣告歸零）};

% 層 3
\node[v3, below=10mm of input, yshift=-22mm] (l3) {%
  \textbf{層 \textcircled{\scriptsize 3}} \quad 子模型確定性匹配\\
  截止日期 $\,\cdot\,$ 能力邊界 $\,\cdot\,$ 拒絕前綴};
\path (l3.west) -- node[layer, minimum width=88mm, draw=none, fill=none] {} (l3.east);

% 層 4
\node[v3e, below=4mm of l3] (l4) {%
  \textbf{層 \textcircled{\scriptsize 4}} \quad 行為向量擴展\\
  拒絕梯度 (8) $\,\cdot\,$ 格式偏好 (3) $\,\cdot\,$ 數值不確定性 (1)};

% 判決融合
\node[fusion, below=4mm of l4] (fuse) {%
  \textbf{判決融合}\\
  7 種互斥狀態 + 置信區間};

% 輸出
\node[output, below=4mm of fuse] (out) {%
  IdentityAssessment \{ verdict, riskFlags, evidence, confidence \}};

% 連線
\draw[arrow] (input.south) -- ++(0,-2mm) -| (l1.north);
\draw[arrow] (input.south) -- ++(0,-2mm) -| (l2.north);
\draw[arrow] (l1.south) -- ++(0,-2mm) -| ([xshift=-22mm]l3.north);
\draw[arrow] (l2.south) -- ++(0,-2mm) -| ([xshift=22mm]l3.north);
\draw[arrow] (l3.south) -- (l4.north);
\draw[arrow] (l4.south) -- (fuse.north);
\draw[arrow] (fuse.south) -- (out.north);

% 側註
\node[label, right=2mm of l1, anchor=west, text width=24mm] {易受自我宣告污染};
\node[label, left=2mm of l2, anchor=east, text width=24mm] {對提示層攻擊穩健};
\node[label, right=2mm of l3, anchor=west, text width=24mm] {可被精準針對};
\node[label, right=2mm of l4, anchor=west, text width=24mm] {政策固化、難以偽造};

\end{tikzpicture}%
}
\caption{四層獨立指紋偵測架構（\S\ref{sec:method}）。每一層使用獨立的訊號通道；
攻擊者必須跨多個通道同時偽造一致行為才能逃過偵測。判決融合產出七種互斥狀態之一,
並附以置信區間。}
\label{fig:architecture}
\end{figure*}
